\begin{abstract}
Reinforcement Learning (RL), specifically Group Relative Policy Optimization (GRPO), has emerged as a transformative technique for unlocking deep reasoning in Large Language Models (LLMs). By incentivizing multi-step Chain-of-Thought (CoT) behaviors, GRPO allows models to transcend simple pattern matching. However, current implementations are plagued by critical bottlenecks: advantage vanishing in uniform groups, systematic correction bias during reflection, and inefficient utilization of noisy, high-entropy error samples. We propose SB-GRPO (Semantic-Balanced GRPO), an advanced framework that integrates dynamic multi-layer reflection with semantic error profiling. Our architecture introduces a differentiable Heuristic Confidence Gate that smoothly transitions between standard and negative-enhanced advantage calculation, preventing gradient stalls. We further employ K-Means clustering on reasoning embeddings to identify systematic logical failure modes and construct Static Balanced Batches to stabilize policy gradients. Finally, we derive an information-theoretic loss scaling factor based on the Shannon Entropy of semantic clusters to prioritize fixing systematic errors over random noise. Evaluated under strict parameter-efficient fine-tuning constraints (4-bit QLoRA) using a 7B model, results show that SB-GRPO outperforms vanilla GRPO by a huge margin of 26.40\% absolute improvement on the MATH-500 benchmark, reaching a highly competitive pass@1 accuracy of 74.80\%. We took a step further when comparing our algorithm with M-GRPO, one of the latest and most prominent multi-layer baselines, demonstrating an additional absolute gain of 2.20\% overall, while significantly outperforming it in complex domains such as Precalculus (+5.7\%) and Intermediate Algebra (+4.9\%). Furthermore, it exhibits substantial gains in elite logical reasoning, effectively tripling the baseline performance to achieve 10.00\% accuracy on the highly challenging AIME 2024 benchmark.

\keywords{Reinforcement Learning \and Large Language Models \and Mathematical Reasoning \and Semantic Entropy}
\end{abstract}